% GENERATED FILE. Edit canonical lessons, then run npm run generate:books.
% canonical-source-hash: fnv1a64:bc6adb9fffe3a89a
% canonical-lessons: MW-C236-gadho, MW-C236-kukro, MW-C236-un, MW-C236-sig, MW-C236-khur

\chapter{Farm Animals}
\label{ch:mw-farm-animals}
\hlchaptermodality{236}

\begin{chapteropening}
\textbf{By the end of this chapter:} \emph{I can name a donkey and a rooster, and wool, horns and hooves.}

Complete the last lesson of chapter 236: I can name a donkey and a rooster, and wool, horns and hooves.
\end{chapteropening}

\section[\texorpdfstring{\mw{गधो} (gadho)}{गधो (gadho)}]{\mw{गधो} (gadho) — a donkey}

\label{lesson:MW-C236-gadho}



\begin{quote}
Before the new one: say the Marwadi for to rear (animals), then the Marwadi for chaff.
\end{quote}

\subsection*{You'll want to know: \mw{गधो}}
\textbf{\mw{गधो}} — \emph{gadho} — \textquotedblleft{}a donkey\textquotedblright{}.

\begin{quote}
\mw{गधो} \mw{भारी} \mw{बोझ} \mw{उठावै} \mw{है।}
\end{quote}

\begin{quote}
\emph{gadho bhārī bojh uṭhāvai hai.}
\end{quote}

\begin{quote}
The donkey carries a heavy load.
\end{quote}

A masculine noun.

\begin{cousinweb}
\textbf{\mw{गधो}} (\emph{gadho}) is inherited from Sanskrit \emph{gardabha}, \textquotedblleft{}an ass\textquotedblright{}; CDIAL lists Hindi \emph{gadahā} and a Gujarati cousin. CDIAL gives no Marwari form in this entry (CDIAL 4054).
\end{cousinweb}

\begin{grammarlens}[title={What you've built}]
One more word for this chapter.
\end{grammarlens}

\subsection*{Your turn}
\begin{itemize}
\raggedright
  \item \emph{Say it:} \emph{gadho}
  \item \emph{Say it:} \emph{gadho}, once more
  \item \emph{Say it:} \emph{gadho bhārī bojh uṭhāvai hai.} — The donkey carries a heavy load.
  \item \emph{Recall:} say the Marwadi for to rear (animals), then the Marwadi for chaff, then say \emph{gadho} again
\end{itemize}

\begin{tcolorbox}[breakable,colback=teal!4,colframe=teal!35!black,title={Before you move on}]
What does \mw{गधो} mean? (\textquotedblleft{}A donkey\textquotedblright{}.) Say it once more.
\end{tcolorbox}
\section[\texorpdfstring{\mw{कूकड़ो} (kūkṛo)}{कूकड़ो (kūkṛo)}]{\mw{कूकड़ो} (kūkṛo) — a rooster}

\label{lesson:MW-C236-kukro}



\begin{quote}
Before the new one: say the Marwadi for chaff, then the Marwadi for a donkey.
\end{quote}

\subsection*{You'll want to know: \mw{कूकड़ो}}
\textbf{\mw{कूकड़ो}} — \emph{kūkṛo} — \textquotedblleft{}a rooster\textquotedblright{}.

\begin{quote}
\mw{कूकड़ो} \mw{तड़कै} \mw{बोलै} \mw{है।}
\end{quote}

\begin{quote}
\emph{kūkṛo taṛkai bolai hai.}
\end{quote}

\begin{quote}
The rooster crows at daybreak.
\end{quote}

A masculine noun.

\begin{cousinweb}
\textbf{\mw{कूकड़ो}} (\emph{kūkṛo}) is inherited from Sanskrit \emph{kukkuṭa}, \textquotedblleft{}a cock\textquotedblright{}, a word CDIAL takes to imitate its call; CDIAL lists Hindi \emph{kukṛā}. CDIAL gives no Marwari form in this entry (CDIAL 3208).
\end{cousinweb}

\begin{grammarlens}[title={What you've built}]
One more word for this chapter.
\end{grammarlens}

\subsection*{Your turn}
\begin{itemize}
\raggedright
  \item \emph{Say it:} \emph{kūkṛo}
  \item \emph{Say it:} \emph{kūkṛo}, once more
  \item \emph{Say it:} \emph{kūkṛo taṛkai bolai hai.} — The rooster crows at daybreak.
  \item \emph{Recall:} say the Marwadi for chaff, then the Marwadi for a donkey, then say \emph{kūkṛo} again
\end{itemize}

\begin{tcolorbox}[breakable,colback=teal!4,colframe=teal!35!black,title={Before you move on}]
What does \mw{कूकड़ो} mean? (\textquotedblleft{}A rooster\textquotedblright{}.) Say it once more.
\end{tcolorbox}
\section[\texorpdfstring{\mw{ऊन} (ūn)}{ऊन (ūn)}]{\mw{ऊन} (ūn) — wool}

\label{lesson:MW-C236-un}



\begin{quote}
Before the new one: say the Marwadi for a donkey, then the Marwadi for a rooster.
\end{quote}

\subsection*{You'll want to know: \mw{ऊन}}
\textbf{\mw{ऊन}} — \emph{ūn} — \textquotedblleft{}wool\textquotedblright{}.

\begin{quote}
\mw{भेड़} \mw{ऊन} \mw{देवै} \mw{है।}
\end{quote}

\begin{quote}
\emph{bheṛ ūn devai hai.}
\end{quote}

\begin{quote}
Sheep give wool.
\end{quote}

A feminine noun.

\begin{cousinweb}
\textbf{\mw{ऊन}} (\emph{ūn}) is inherited from Sanskrit \emph{ūrṇā}, \textquotedblleft{}wool\textquotedblright{}; CDIAL lists Hindi and Gujarati \emph{ūn}. CDIAL gives no Marwari form in this entry (CDIAL 2424).
\end{cousinweb}

\begin{grammarlens}[title={What you've built}]
One more word for this chapter.
\end{grammarlens}

\subsection*{Your turn}
\begin{itemize}
\raggedright
  \item \emph{Say it:} \emph{ūn}
  \item \emph{Say it:} \emph{ūn}, once more
  \item \emph{Say it:} \emph{bheṛ ūn devai hai.} — Sheep give wool.
  \item \emph{Recall:} say the Marwadi for a donkey, then the Marwadi for a rooster, then say \emph{ūn} again
\end{itemize}

\begin{tcolorbox}[breakable,colback=teal!4,colframe=teal!35!black,title={Before you move on}]
What does \mw{ऊन} mean? (\textquotedblleft{}Wool\textquotedblright{}.) Say it once more.
\end{tcolorbox}
\section[\texorpdfstring{\mw{सींग} (sī̃g)}{सींग (sī̃g)}]{\mw{सींग} (sī̃g) — a horn}

\label{lesson:MW-C236-sig}



\begin{quote}
Before the new one: say the Marwadi for a rooster, then the Marwadi for wool.
\end{quote}

\subsection*{You'll want to know: \mw{सींग}}
\textbf{\mw{सींग}} — \emph{sī̃g} — \textquotedblleft{}a horn\textquotedblright{}.

\begin{quote}
\mw{हिरण} \mw{माथै} \mw{सींग} \mw{है।}
\end{quote}

\begin{quote}
\emph{hiraṇ māthai sī̃g hai.}
\end{quote}

\begin{quote}
The deer has horns on its head.
\end{quote}

A masculine noun.

\begin{cousinweb}
\textbf{\mw{सींग}} (\emph{sī̃g}) is inherited from Sanskrit \emph{śṛṅga}, \textquotedblleft{}a horn\textquotedblright{}; CDIAL lists Hindi \emph{sī̃g}, \textquotedblleft{}a horn\textquotedblright{}, and a Gujarati cousin. CDIAL gives no Marwari form in this entry (CDIAL 12583).
\end{cousinweb}

\begin{grammarlens}[title={What you've built}]
One more word for this chapter.
\end{grammarlens}

\subsection*{Your turn}
\begin{itemize}
\raggedright
  \item \emph{Say it:} \emph{sī̃g}
  \item \emph{Say it:} \emph{sī̃g}, once more
  \item \emph{Say it:} \emph{hiraṇ māthai sī̃g hai.} — The deer has horns on its head.
  \item \emph{Recall:} say the Marwadi for a rooster, then the Marwadi for wool, then say \emph{sī̃g} again
\end{itemize}

\begin{tcolorbox}[breakable,colback=teal!4,colframe=teal!35!black,title={Before you move on}]
What does \mw{सींग} mean? (\textquotedblleft{}A horn\textquotedblright{}.) Say it once more.
\end{tcolorbox}
\section[\texorpdfstring{\mw{खुर} (khur)}{खुर (khur)}]{\mw{खुर} (khur) — a hoof}

\label{lesson:MW-C236-khur}



\begin{quote}
Before the new one: say the Marwadi for wool, then the Marwadi for a horn.
\end{quote}

\subsection*{You'll want to know: \mw{खुर}}
\textbf{\mw{खुर}} — \emph{khur} — \textquotedblleft{}a hoof\textquotedblright{}.

\begin{quote}
\mw{खुर} \mw{कड़क} \mw{है।}
\end{quote}

\begin{quote}
\emph{khur kaṛak hai.}
\end{quote}

\begin{quote}
The hoof is hard.
\end{quote}

A masculine noun.

\begin{cousinweb}
\textbf{\mw{खुर}} (\emph{khur}) is inherited from Sanskrit \emph{khura}, \textquotedblleft{}a hoof\textquotedblright{}; CDIAL lists Hindi and Gujarati \emph{khur}, \textquotedblleft{}a hoof\textquotedblright{}. CDIAL gives no Marwari form in this entry (CDIAL 3906).
\end{cousinweb}

\begin{grammarlens}[title={What you've built}]
One more word for this chapter.
\end{grammarlens}

\subsection*{Your turn}
\begin{itemize}
\raggedright
  \item \emph{Say it:} \emph{khur}
  \item \emph{Say it:} \emph{khur}, once more
  \item \emph{Say it:} \emph{khur kaṛak hai.} — The hoof is hard.
  \item \emph{Recall:} say the Marwadi for wool, then the Marwadi for a horn, then say \emph{khur} again
\end{itemize}

\begin{tcolorbox}[breakable,colback=teal!4,colframe=teal!35!black,title={Before you move on}]
What does \mw{खुर} mean? (\textquotedblleft{}A hoof\textquotedblright{}.) Say it once more.
\end{tcolorbox}
